\hypertarget{outer.extension}{}
\begin{proposition}[Odd-power extension]\label{prop:outer-extension}
Let a primitive $F$ in \eqref{eq:lienard} have a finite nonempty
collection of nested hyperbolic cycles. For every odd $p>\deg F$, there
are arbitrarily small nonzero coefficients of $x^p$ for which the
perturbed system retains these cycles and has a further hyperbolic
cycle outside them. In particular, for each $k\ge3$ the support
$\{1,9,10\}\cup\{2i+3:4\le i\le k\}$ admits an open set with at least
$k$ hyperbolic cycles.
\end{proposition}
\begin{proof}
The outer trapping step follows the higher-damping construction of
\cite[Remark~1.2]{DPR} and the odd-primitive extension in
\cite[\S3.2]{DMH}. We show that adjoining any odd power $x^p$ with $p>\deg F$ to a
polynomial having a finite nonempty hyperbolic nest gives one more cycle.
Take the outermost selected cycle and, by reversing time and replacing
$y$ by $-y$ if necessary, make it repelling. In a thin exterior collar
choose a smooth Jordan curve $C_i$ crossed strictly outwards. Such curves
are small constant transverse-coordinate curves in Floquet coordinates
at a repelling cycle. Choose the collar so that $C_i$ meets the upper
axis just once, at height $\beta_i$.

For $P_\delta(x)=F(x)+\delta x^p$, all selected cycles and the strict
crossing persist for sufficiently small $\delta>0$. Fix one such
$\delta_0$ and put $P=P_{\delta_0}$. Define
\[
 v(x)=\int_0^x\frac{\dd s}{1+P(s)^2},\qquad
 W=\tfrac12(x^2+y^2)-yv(x).
\]
The function $v$ is bounded, $|v(x)|<|x|$ for $x\ne0$, and
$0<v'\le1$. Completing the square gives
$2W=(y-v)^2+x^2-v^2$. The last two terms increase strictly with $|x|$
on each half-line. Large levels are therefore smooth ovals, each meeting
the positive axis once. Direct differentiation gives
\[
 \dot W=-x(P-v)-v'(y-P/2)^2+v'P^2/4 .
\]
The last term is at most $1/4$. The first tends to $-\infty$ as
$|x|\to\infty$ because the leading coefficient of $P$ is positive and
$p$ is odd. For bounded $x$, the squared term gives the same conclusion
as $|y|\to\infty$. Choose an enclosing level $C_o$ with $\dot W<0$ and
upper-axis height $\beta_o$. Keep both ovals fixed while varying $\delta$
in a small open interval about $\delta_0$. The annulus between the ovals
is positively invariant and contains no equilibrium.

For each $\beta\in[\beta_i,\beta_o]$ the forward orbit returns to the
positive axis. Indeed its omega-limit set in the annulus is a periodic
orbit by Poincar\'e--Bendixson, and a flow box at that orbit's transverse
axis crossing forces a return. Transversality and the polynomial vector
field give real-analytic dependence of the full return $\Pi_\delta$ on
the launch height. At the boundaries $\Pi_\delta(\beta_i)>\beta_i$ and
$\Pi_\delta(\beta_o)<\beta_o$. Thus for the fixed $\delta_0$ at least
one exterior fixed point is isolated and of odd multiplicity.
Hyperbolicity requires a simple fixed point. We now vary $\delta$
inside the interval on which both crossings persist.
At any zero of $\Pi_\delta(\beta)-\beta$ both reflected half-returns exist
in a neighborhood. Differentiating
$b_-(\Pi_\delta(\beta),\delta)=b_+(\beta,\delta)$ there gives at the zero
\[
 \partial_\delta\Pi_\delta=
 \frac{\partial_\delta b_+-\partial_\delta b_-}{\partial_\beta b_-}<0
\]
by Lemma~\ref{lem:variational}(ii). The local zero curves therefore
project to the $\delta$-axis with critical values of measure zero
(Sard's theorem, as in Lemma~\ref{lem:hyperbolicity}). For a regular
$\delta$ the opposite boundary signs give a hyperbolic cycle in the
annulus. The original nest persists inside $C_i$, and the new coefficient
is nonzero.

Alternatively, if $x$ is already in the support, hold $\delta_0$ fixed
and vary the linear coefficient $a_1$. The same annulus persists.
In \eqref{eq:b-param} the two variations are $+r$ and $-r$,
so $\partial_{a_1}(b_+-b_-)>0$ at every return zero. Equivalently
$\det((y-F,-x),(-x,0))=-x^2$ along the orbit, strictly negative
off the axis. The local zero curves and the same regular-value argument
give a simple exterior cycle while $a_1$ stays nonzero. Undoing the time and ordinate reversal, when used, restores
the original coefficients of $F$.

For the stated support take $d=10$ in Section~3 of the electronic
supplement. A regular value of the trinomial matching curve gives
three simple returns of $C(x^{10}+x^9-ax)$ with $C,a>0$. The three
hyperbolic cycles enclose the unique equilibrium, so they form the
nest required above. This is the degree-ten case originally obtained
in \cite[proof of Theorem~1.2]{CJLZ}. Successively adjoining
$11,13,15,\ldots$ as exponents preserves the earlier cycles and adds
one each time. Persistence at each finite stage gives the open set.
\end{proof}
